\documentclass[../../../include/open-logic-section]{subfiles}

\begin{document}

\olfileid{sth}{card-arithmetic}{fix}
\olsection{نقاط $\aleph$ الثابتة}

اقترحنا في \olref[spine][]{chap} أن بديهية الاستبدال بحاجة إلى تبرير،
لأنها تجبر التراتبية على أن تكون بالغة الارتفاع. وبعد أن درسنا شيئاً من
حساب الكاردينالات، نستطيع أن نقدم مثالاً صغيراً على ارتفاع التراتبية.

من الواضح أن $0 < \aleph_0$ و$1 < \aleph_1$ و$2 < \aleph_2$، وهكذا؛
بل إن فرق الحجم لا يزداد إلا \emph{كبراً} مع كل خطوة. ولذلك يغري هذا
بالتخمين أن $\kappa< \aleph_\kappa$ لكل عدد ترتيبي $\kappa$.

لكن هذا التخمين \emph{كاذب} في $\ZFC$. فنستطيع في الواقع أن نثبت وجود
\emph{نقاط ثابتة لـ$\aleph$}، أي كاردينالات $\kappa$ تحقق
$\kappa=\aleph_\kappa$.

\begin{prop}\ollabel{alephfixed}
توجد نقطة ثابتة لـ$\aleph$.
\end{prop}

\begin{proof}
عرّف بالعود:
\begin{align*}
	\kappa_0 &= 0\\
	\kappa_{n+1} &= \aleph_{\kappa_n}\\
	\kappa&= \bigcup_{n < \omega}\kappa_n
\end{align*}
يكون $\kappa$ كاردينالاً بحسب
\olref[cardinals][classing]{unioncardinalscardinal}. والآن:
\[
	\kappa= \bigcup_{n < \omega} \kappa_{n+1} =
	\bigcup_{n < \omega}\aleph_{\kappa_n} =
	\bigcup_{\alpha < \kappa}\aleph_\alpha = \aleph_\kappa
\]
\end{proof}

كتب بولوس مرة مقالة عن نقطة $\aleph$ الثابتة نفسها التي أنشأناها للتو.
وبعد أن لاحظ وجود $\kappa$ في مطلع مقالته قال:
\begin{quote}
	[$\kappa$] عدد \emph{كبير جداً} بحسب معايير من لم يسبق لهم التعرض
	لنظرية المجموعات؛ وهو، فيما يبدو لي، كبير إلى حد يجعل صدق أي نظرية
	تتضمن بين أقوالها زعم وجود $\kappa$ من الأشياء على الأقل موضع شك.
	\citep[p.~257]{Boolos2000}
\end{quote}
ثم ختم مقالته أخيراً بالسؤال:
\begin{quote}
	[هل] نشك في أنه، أياً يكن الحال في بداية القصة، حين نكون قد بلغنا هذا
	الموضع تكون العجلات قد أخذت تدور في فراغ، ولم نعد نصغي إلى وصف شيء
	واقع؟
	\citep[p.~268]{Boolos2000}
\end{quote}
فإذا كنا قد تجاوزنا بالفعل «أي شيء واقع»، وجب أن نوجه اللوم مباشرة إلى
الاستبدال، لأن هذه البديهية هي التي تتيح نجاح برهاننا. وعندئذ، فيما
نفترض، سيحتاج بولوس إلى مراجعة قوله، قبل بضعة عقود، إن الاستبدال لا تترتب
عليه آثار «غير مرغوب فيها» (انظر
\olref[replacement][extrinsic]{sec}).

ولكن هل وجود $\kappa$ سيئ إلى هذا الحد؟ قد يفيد هنا أن ننظر في
\emph{مفارقة تريسترام شاندي} عند راسل. يوثق تريسترام شاندي حياته في
يومياته، لكن تسجيل يوم واحد يستغرق منه سنة. ومع مرور كل سنة يزداد
تريسترام تأخراً: فبعد سنة لا يكون قد سجل إلا يوماً واحداً، ويكون قد عاش
364 يوماً غير مسجلة؛ وبعد سنتين لا يكون قد سجل إلا يومين، ويكون قد عاش
728 يوماً غير مسجلة؛ وبعد ثلاث سنوات لا يكون قد سجل إلا ثلاثة أيام،
ويكون قد عاش 1092 يوماً غير مسجلة
\dots\footnote{نتجاهل السنوات الكبيسة.} ومع ذلك، إذا كان تريسترام
\emph{خالداً} فسيتمكن من تسجيل كل يوم، لأنه سيسجل اليوم $n$ في السنة
$n$ من حياته. ولذلك ستكون لتريسترام، «في نهاية الزمان»، يوميات كاملة.

فلماذا يختلف هذا إلى ذلك الحد عن فكرة أن $\alpha$ أصغر من
$\aleph_\alpha$---بل أصغر منه على نحو متزايد ويائس---حتى نصل إلى
$\kappa$، حيث نلحق به ويصبح $\kappa = \aleph_\kappa$؟

فلندع ذلك جانباً، ولنفترض قبولنا $\ZFC$، ولنختتم بمزيد من المتعة بشأن
إنشاء النقاط الثابتة. تثبت النتائج الثلاث الآتية، على نحو حدسي، أن ثمة
نقطة غير تافهة تكون عندها التراتبية عريضة بقدر ارتفاعها:

\begin{prop}\ollabel{bethfixed}
توجد نقطة ثابتة لـ$\beth$، أي يوجد $\kappa$ بحيث
$\kappa= \beth_\kappa$.
\end{prop}

\begin{proof}
كما في \olref{alephfixed}، باستعمال «$\beth$» بدلاً من «$\aleph$».
\end{proof}

\begin{prop}\ollabel{stagesize}
$\card{V_{\omega+\alpha}} = \beth_{\alpha}$. وإذا كان
$\omega \ordtimes \omega \leq \alpha$، فإن
$\card{V_\alpha} = \beth_\alpha$.
\end{prop}

\begin{proof}
يثبت الادعاء الأول باستقراء بسيط ما فوق المتناهي. ويتبع الادعاء الثاني،
لأنه إذا كان $\omega \ordtimes \omega \leq \alpha$ فإن
$\omega + \alpha = \alpha$. ولإثبات ذلك نستعمل حقائق عن حساب الأعداد
الترتيبية من \olref[ord-arithmetic][]{chap}. لاحظ أولاً أن
$\omega \ordtimes \omega = \omega \ordtimes (1 \ordplus \omega) =
(\omega \ordtimes 1) \ordplus (\omega\ordtimes\omega) = \omega
\ordplus (\omega \ordtimes \omega)$. فإذا كان
$\omega \ordtimes \omega \leq \alpha$، أي كان
$\alpha = (\omega\ordtimes\omega) \ordplus \beta$ لبعض $\beta$، فإن
$\omega \ordplus \alpha = \omega \ordplus
((\omega \ordtimes \omega) \ordplus \beta) =
(\omega \ordplus (\omega \ordtimes \omega)) \ordplus \beta =
(\omega \ordtimes \omega) \ordplus \beta = \alpha$.
\end{proof}

\begin{cor}
يوجد $\kappa$ بحيث $\card{V_\kappa} = \kappa$.
\end{cor}

\begin{proof}
ليكن $\kappa$ نقطة ثابتة لـ$\beth$ كما تضمنها
\olref{bethfixed}. ومن الواضح أن
$\omega \ordtimes \omega < \kappa$. ومن ثم، بحسب
\olref{stagesize}، $\card{V_\kappa} = \beth_\kappa= \kappa$.
\end{proof}

تحت $V_\kappa$ من المراحل بقدر ما في $V_\kappa$ من !!{element}.
إذن، على نحو حدسي، يكون $V_\kappa$ عريضاً بقدر ارتفاعه. وهذا شبيه جداً
بتريسترام شاندي: ننتقل من مرحلة إلى تليها بأخذ \emph{مجموعات القوى}،
فنجعل تراتبيتنا \emph{أكبر بكثير} في كل خطوة. لكن «في النهاية»، أي عند
المرحلة $\kappa$، يلحق عرض التراتبية بارتفاعها.

قد يسأل المرء: \emph{كم مرة يساوي عرض التراتبية ارتفاعها؟} والجواب:
\emph{بعدد الأعداد الترتيبية.} لكن هذا يحتاج إلى شيء من الشرح.

نعرّف حداً $\tau$ كما يأتي. لأي $A$، ضع:
\begin{align*}
	\tau_0(A) & \defis \cardsucc{\card{A}}\\
	\tau_{n+1}(A) & \defis \beth_{\tau_n(A)}\\
	\tau(A) & \defis \bigcup_{n < \omega}\tau_n(A)
\intertext{كما في \olref{bethfixed}، تكون $\tau(A)$ نقطة ثابتة لـ$\beth$
لكل $A$، ومن تعريف $\tau_0(A)$ يكون
$\card{A} < \tau(A)$ بوضوح. فانظر الآن إلى التعريف العودي الآتي:}
	W_0 &\defis 0\\
	W_{\alpha + 1} & \defis \tau(W_\alpha)\\
	W_\alpha & \defis \bigcup_{\beta < \alpha} W_\beta
	\text{، إذا كان $\alpha$ حدياً.}
\end{align*}
هذا البناء معرّف لجميع الأعداد الترتيبية. وكل $W_\alpha$ كاردينال؛ وفي
كل خطوة خلفية لدينا
$W_\alpha=\card{W_\alpha}<\tau(W_\alpha)=W_{\alpha+1}$، بينما نأخذ
الاتحاد في المراحل الحدية. لذلك تكون الخريطة
$\alpha\mapsto W_\alpha$ متزايدة قطعاً، ومن ثم فهي، على نحو حدسي،
«!!a{injection}» من الأعداد الترتيبية إلى نقاط $\beth$ الثابتة. وتماماً
كما سبق، يكون $V_{W_\alpha}$ عريضاً بقدر ارتفاعه، لكل $\alpha$.

\end{document}
